\documentclass[11pt,a4paper]{article}

\usepackage[utf8]{inputenc}
\usepackage[T1]{fontenc}
\usepackage[margin=2.4cm]{geometry}
\usepackage{parskip}
\usepackage{titlesec}
\usepackage{enumitem}
\usepackage{booktabs}
\usepackage{xcolor}
\usepackage{tcolorbox}
\usepackage{hyperref}
\usepackage{fancyhdr}

\definecolor{accent}{RGB}{30,80,150}
\definecolor{softgray}{RGB}{245,246,248}

\titleformat{\section}{\Large\bfseries\color{accent}}{\thesection}{0.8em}{}
\titleformat{\subsection}{\large\bfseries}{\thesubsection}{0.6em}{}

\hypersetup{
  colorlinks=true,
  linkcolor=accent,
  urlcolor=accent,
  pdftitle={SkillOpt for AgentIA --- Proposal Brief},
  pdfauthor={AgentIA Engineering}
}

\pagestyle{fancy}
\fancyhf{}
\fancyhead[L]{\small\color{accent}AgentIA $\cdot$ Skill Layer Proposal}
\fancyhead[R]{\small\thepage}
\renewcommand{\headrulewidth}{0.4pt}

\newtcolorbox{keybox}[1][]{
  colback=softgray, colframe=accent, boxrule=0.4pt,
  arc=2pt, left=8pt, right=8pt, top=6pt, bottom=6pt, #1
}

\title{\vspace{-2cm}\textbf{SkillOpt for AgentIA}\\[0.2em]
\large A bounded text-space optimization layer for procedural reliability}
\author{AgentIA Engineering}
\date{\today}

\begin{document}
\maketitle
\thispagestyle{fancy}

\begin{keybox}
\textbf{In one paragraph.} AgentIA generates Java\,21 / Spring\,Boot\,3
microservice scaffolding and validates it with \texttt{mvn test -o}. Its
recurring failures are procedural: the model \emph{knows} a rule (use
\texttt{jakarta.*}, not \texttt{javax.*}; controllers must not inject
repositories directly) but forgets to apply it in a given session. We
propose adding a \emph{skill layer} --- a small set of triggerable
procedural documents prepended to the generation prompt --- and, once
its value is demonstrated, an optimizer that improves those documents
under a held-out validation gate. The optimizer is SkillOpt
(Microsoft Research, 2026). The plan is staged: \emph{inject and
measure} first (1--2 weeks); \emph{optimize} only if the measurement
shows signal (4 weeks).
\end{keybox}

\section{The problem}

AgentIA's generation pipeline ends in \texttt{VERIFIED} or
\texttt{BLOCKED}. A \texttt{BLOCKED} session has exhausted the bounded
repair loop (max 3 attempts) and requires human intervention. Two
distinct failure classes produce this outcome:

\begin{description}[leftmargin=1.2em]
  \item[Capability gap.] The model does not know the rule. A better
  reminder will not help, because there is nothing to remind.
  \item[Application gap.] The model knows the rule but fails to apply
  it in a specific session. A persistent, triggered reminder can help.
\end{description}

The SkillOpt layer targets the \emph{second} class. This distinction
matters for scope: the mechanism improves reliability within what
AgentIA already covers, not coverage of what it does not (Spring
Security configuration, reactive programming, Saga/CQRS, the Azure
stack). Coverage expansion is a separate investment.

\section{What a skill is}

A skill is \emph{not} a cheat sheet (static facts). It is a
\emph{procedural policy} with three parts:

\begin{enumerate}[leftmargin=1.4em]
  \item \textbf{Trigger} (\texttt{when\_to\_apply}) --- the situation in
  which the skill becomes relevant.
  \item \textbf{Procedure} (\texttt{rules}) --- an ordered set of
  actions to take once the trigger fires.
  \item \textbf{Granularity} --- task-level (applies to a whole task) or
  event-driven (fires on a specific execution signal).
\end{enumerate}

Representative skill (adapted from CODESKILL's Java/Maven examples):

\begin{keybox}
\textbf{Title.} Handle missing Maven test-jar dependencies.\\
\textbf{When to apply.} A Maven build fails with a dependency-resolution
error for a test-jar that has not yet been produced.\\
\textbf{Rules.}
(1) Inspect the failing module's POM to confirm dependency type and scope.
(2) Check the producing module's POM for an existing plugin execution.
(3) If the test jar is only generated when tests run, run tests or
invoke the test-jar goal explicitly.
(4) Do not retry with tests skipped --- skipping prevents the artifact.
(5) Run a full install so the artifact lands in the local repository.
(6) Verify presence before retrying the original command.
\end{keybox}

At deployment, skills are \emph{prepended} to the generation prompt.
The target model remains frozen. No inference-time overhead beyond the
prompt tokens.

\section{What SkillOpt is}

SkillOpt treats the skill document as a \textbf{trainable parameter}
and builds a training loop around it --- in text space, with the same
discipline that makes weight-space optimization reproducible.

\subsection{The four controls}

\begin{description}[leftmargin=1.2em]
  \item[Bounded edit budget (learning rate).] Each step applies at most
  \(L_t\) atomic edits (default \(4\), cosine decay, floor \(2\)).
  Prevents unbounded rewriting; makes each step small enough that its
  effect is measurable.
  \item[Held-out validation gate (trust region).] Every candidate
  document is scored on a \emph{selection split} --- sessions never
  used to motivate the edits. Accepted only if the score
  \emph{strictly} improves. Ties rejected. This is what makes the loop
  auditable and non-degrading.
  \item[Rejected-edit buffer (negative feedback).] Rejected edits and
  their score drops are stored and fed back to the next reflection
  call, so the optimizer avoids repeating failed directions.
  \item[Epoch-wise slow/meta update (momentum).] At each epoch
  boundary, the same 20 tasks are re-run under the previous and
  current documents; the optimizer writes a \emph{longitudinal
  guidance block} into a protected section that step-level edits
  cannot touch. Catches drift that adjacent-step reflection cannot see.
\end{description}

\subsection{The loop}

\begin{keybox}
\textbf{Inner loop (per step).}\\
rollout on training set $\to$ split into success/failure minibatches
$\to$ optimizer proposes edits $\to$ hierarchical merge, dedupe, rank
$\to$ clip to top \(L_t\) $\to$ apply $\to$ candidate $\to$
\textsc{Gate}: accept iff selection score strictly improves, else
reject and log.

\medskip
\textbf{Outer loop (per epoch, 4 total).}\\
sample 20 tasks $\to$ re-run under previous and current skill $\to$
categorize regressions / persistent failures / improvements / stable
successes $\to$ optimizer writes longitudinal guidance $\to$ write to
protected section $\to$ same gate.
\end{keybox}

\subsection{What it produces}

A single \texttt{best\_skill.md}, roughly \(300\)--\(2{,}000\) tokens,
fully inspectable. No inference-time model calls; the document is
prepended to the prompt and nothing else changes at deployment.

\section{The evidence}

The paper's evaluation is unusually broad: six benchmarks, seven target
models, three execution harnesses. Headline results:

\begin{itemize}[leftmargin=1.4em]
  \item Best or tied-best on \textbf{52 of 52} evaluated
  (model, benchmark, harness) cells.
  \item Average gain on GPT-5.5 direct chat: \(+\,23.5\) points over
  no-skill; \(+\,24.8\) under Codex; \(+\,19.1\) under Claude Code.
  \item Beats an \emph{oracle} baseline that picks the best of six
  competing methods per cell by \(+\,5.4\) points.
  \item \textbf{Edit economy:} \(1\)--\(4\) accepted edits (median
  \(2.5\)) produce \(+\,29\) to \(+\,39\) points on procedural
  benchmarks. OfficeQA's \(+\,39.0\) arises from a single accepted edit.
  \item \textbf{Cross-harness transfer:} a skill trained under Codex
  transfers to Claude Code with \(+\,59.7\) over baseline --- exceeding
  the in-domain result.
\end{itemize}

\subsection{The ablation that matters most}

Removing both the meta-skill and the slow update drops SpreadsheetBench
from \(77.5\) to \(55.0\) --- a \(22.5\) point degradation, the largest
in the paper. The epoch-level longitudinal view is load-bearing, not
decorative.

\section{Honest boundaries}

\begin{itemize}[leftmargin=1.4em]
  \item \textbf{No code benchmark.} None of the six benchmarks is code.
  The mechanism is domain-agnostic (it operates on text plus a
  verifier), but the numbers do not transfer. AgentIA's pilot would be
  the first application to Java generation.
  \item \textbf{Requires an automated verifier.} The paper's Limitations
  section is explicit. \texttt{mvn test -o} is exactly this.
  \item \textbf{One document, not a bank.} SkillOpt optimizes a single
  skill artifact. Bank-level curation (which skills to keep or retire)
  is a different problem (SkillBrew, Ratchet) deferred to later specs.
  \item \textbf{No weight changes.} This is a feature: no training
  infrastructure, no data pipeline, no fine-tuning.
\end{itemize}

\section{The staged plan}

\subsection{Spec 1 --- inject and measure (1--2 weeks)}

\begin{enumerate}[leftmargin=1.4em]
  \item \textbf{Task 0.} Query \texttt{studio.db} for the historical
  pass/fail distribution and repair-parser error classification. Triage
  the five candidate domains: skill-layer appropriate,
  deterministic-fixer appropriate, or not observed.
  \item Hand-write seed skills for surviving domains.
  \item Add \texttt{skills} table; add \texttt{skill\_injection\_node}
  before \texttt{validator\_node}; log every injection
  (\texttt{session\_id}, \texttt{skill\_ids}, \texttt{versions},
  \texttt{injection\_point}).
  \item Re-inject on every repair attempt.
  \item Run a four-condition A/B pilot (minimum \(15\) sessions per
  condition): no-skills, skills, sham injection (isolates ``longer
  prompt'' from ``this content''), skills \(+\) deterministic fixer.
  \item \textbf{Primary metric:} \texttt{mvn test -o} pass rate.
  \textbf{Secondary:} reviewer tag
  (\texttt{as-is} / \texttt{minor edits} / \texttt{rewrite}).
\end{enumerate}

\subsection{Spec 2 --- optimize (4 weeks, only if Spec 1 shows signal)}

Build the SkillOpt loop: rollout collector, reflector (using the
paper's published prompts), merger, ranker, applier, gate, rejected-edit
buffer, slow/meta update. Optimizer model is a \emph{different} provider
than the target, routed through the existing \texttt{LLMFactory}.

\section{The pilot in detail}

\begin{table}[h]
\centering
\small
\begin{tabular}{@{}lll@{}}
\toprule
Condition & Prompt addition & Purpose \\
\midrule
\texttt{no\_skills} & none & Baseline \\
\texttt{skills} & active skill documents & Treatment \\
\texttt{sham} & \(\sim\)1.5K tokens of irrelevant text & Prompt-length control \\
\texttt{skills\_plus\_fixer} & skills + deterministic post-processor & Fixer interaction \\
\bottomrule
\end{tabular}
\end{table}

\noindent\textbf{Causal design.} Random assignment within the same
agent version and time window, not skill-on-now vs.\ skill-off-historically.
\textbf{Skill versioning.} Freeze skill content for the full pilot
window; version-tag every session.

\section{Decision points}

\begin{enumerate}[leftmargin=1.4em]
  \item \textbf{After Task 0.} Which domains survive triage? Which are
  better handled by deterministic fixers?
  \item \textbf{After Spec 1.} Does injection move the primary metric?
  If no, stop. If yes, build Spec 2.
  \item \textbf{Before Spec 2.} Spec 2 versus coverage expansion
  (Spring Security config generation) --- evaluated on equal footing
  against the original economic framing, not by momentum.
\end{enumerate}

\begin{keybox}
\textbf{The arithmetic worth stating explicitly.} Reliability
improvements on a \(30\)--\(40\%\) covered slice cannot mathematically
reach a target defined against \(100\%\) of the requisition while the
remaining \(60\)--\(70\%\) is uncovered. The reliability work is
necessary infrastructure for any coverage expansion, but it is not a
substitute for it. The resourcing decision after Spec 1 should reflect
that.
\end{keybox}

\section{References}

\begin{itemize}[leftmargin=1.4em]
  \item Li et al. \emph{SkillOpt: A controllable text-space optimizer
  for agent skills.} arXiv:2605.23904, 2026.
  \item Li et al. \emph{CODESKILL: Learning self-evolving skills for
  coding agents.} arXiv:2605.25430, 2026.
  \item Hu et al. \emph{SkillBrew: Multi-objective curation of skill
  banks for LLM agents.} arXiv:2605.29440, 2026.
  \item Eliseeva et al. \emph{EnvBench: A benchmark for automated
  environment setup.} arXiv:2503.14443, 2025.
\end{itemize}

\end{document}